% Section 2: definitions (P4)
Throughout this paper we use the dimensional vocabulary below, and each dimension is chosen to
vary independently: matching on one dimension tells you nothing about the others, and improving
one does not require changing the rest.

\subsection{What reliability is}
Operationally, a configuration is reliable only relative to three named terms: an envelope of failure kinds it must not produce, a budget it must hold, and an assurance path that intervenes as it drifts. ``Reliable'' carries those relative terms; the selection procedure of \cref{sec:selection} is how they are chosen. Reliability is not benchmark accuracy or fluency \cite{measwhatmatters2025,quantconstruct2026}. A system can raise benchmark scores while
becoming less deployable, more confidently wrong, less able to refuse, more
expensive to operate within a budget \cite{leaderillusion2025}. Reliability is a property of the
\emph{system configuration} as well: the contract, thresholds, monitors and escalation paths
below belong to that configuration, not to a single weight checkpoint. That is why our
unit of analysis per \cref{sec:stack} is the assembled stack instance.

Classical reliability engineering supplies the vocabulary this section borrows, and three
distinctions do most of the work. Reliability is the continuity of correct service and
availability is readiness for correct service \cite{avizienis2004}; safety is freedom from
losses to people and environment, a property of the controlled system \cite{leveson2011}. In
deployment terms we read availability as uptime, which can be high while outputs are wrong; that
gap is why we keep the two apart \evjudge{}. Second, reliability is
compositional: a chain of redundant components is only as independent as its failure causes
allow, and common-cause failures defeat redundancy that looks adequate on paper, which is why a
stack that adds a monitor and a verifier has added no independence unless the model, the monitor,
and the verifier fail for different reasons \cite{leveson2011,leveson2018stpa}. Third, accident analysis separates latent
conditions, present in design and organization long before an event, from active failures at the
point of operation \cite{leveson2011}. Our inherited failure envelope is the stack's
latent-condition accounting: the failure a layer defers is a latent condition for the layer
below. Conventional analysis grades each failure mode by severity, occurrence and detectability
(FMECA, IEC~60812 \cite{iec60812}); we use it to rank layers by what they leave unrecoverable;
hazard frequency is a different axis.
Of these terms, only \emph{reliability debt} is a coinage;
the envelope is our adaptation of latent-condition accounting.

Every claim in this paper carries one of six evidence classes: [T] theory under stated assumptions,
[E] independent empirical, [A] author-reported, [V] vendor-reported, [S] verified synthesis,
[W] watch/pending. Four annotations record what a claim is reacting to, while the class carries
the strength of the evidence: [C] contradiction, [N]
named gap, [GAP] unmeasured link, [J] our analytic judgement.
\cref{tab:evidence-classes} is the single alphabet every other surface cites, and
\cref{sec:protocol} gives the full protocol.

\subsection{The dimensions}
\begin{description}[leftmargin=1.5em, itemsep=1pt]
\item[Inherited failure envelope.] The set of failure modes a layer's choice makes unrecoverable
downstream: the choice cannot remove them, it can only relocate them. When a fix moves a failure
to another layer rather than removing it, we call that \emph{reliability debt}.
\item[Well-posedness under distribution shift.] A task is well-posed when its answer space and
input distribution can both be stated; a model is reliable only relative to that pair.
Deployment drifts the second term, which is why shift is a first-class modifier in
\cref{sec:selection}, and why the classifier-baseline literature, with
the longest experience of this instability, anchors the dimension
\cite{id231108298,id250315850}.
\item[Calibration and selective prediction.] Calibration is measured confidence--accuracy
agreement. Selective prediction turns that agreement into a decision by fixing a cutoff with a
stated error rate \cite{id240718418,id260704430}. The formulation is older than the LLM
literature: fix a cut on a confidence score, accept the region it selects, and report the error on
the accepted set \cite{geifman2017selective,geifman2019selectivenet}. What changes in the LLM
setting is the calibration error of the score, confidence that need not be monotone across an
option set, and a rejected sample that has already cost a generation. The direction of
miscalibration (over- vs.\ under-confidence) changes which way a confidence-gated cascade fails
\cite{id260933843}; we therefore track direction as well as magnitude.
\item[Validity versus truth.] Syntactic/schema correctness of an output is separable from
semantic correctness of its content \cite{jsonschemabench2025,xgrammar2024}. Conflating
these is the source of most spurious ``guaranteed'' deployment claims, and the asymmetry
shapes Layer 4 (\cref{sec:contracts}).
\item[Grounding.] For generative output, whether a claim is supported by available evidence.
The trustworthy-RAG literature treats grounding as a collection of techniques organized by
dimension, each opening new failure paths of its own (stale sources, retrieval misses, unfaithful
use) \cite{id250206872,id240910102}.
\item[Error cost and reversibility.] Errors differ by consequence severity and by whether
they can be undone. This becomes the L0 risk budget: the escalation band of \cref{fig:map} is
derived from it, and the band's width is our answer to how much error a task can absorb.
\item[Prediction unit and generation schedule.] What a system predicts --- a token, a continuous
latent, or a discrete code spanning several tokens --- and the order in which predictions
condition on each other are separable from both the training signal and the architecture: the
same objective admits a left-to-right chain, an iterative denoising sweep, or a coarse-to-fine
level order (\cref{fig:schedule-atlas}). The schedule fixes where commitment points fall, and
with them whether early commitments remain revisable.
\item[Latency, cost and budget.] Deployment reliability is bounded by operational budget;
a configuration that is reliable at unit cost can be unreliable at required throughput. We
treat budget as a selection constraint; maximizing it is a different objective.
\item[Execution configuration.] The inference engine, the numeric precision, the cache settings
and the co-resident traffic a model is served under: properties of the deployment; the weights
alone do not carry them. They enter the computed function, so an error rate belongs to the
configuration it was measured on (\cref{sec:l5-bridge}).
\item[Monitoring and drift response.] This dimension is the assurance layer's own reliability:
monitors can be evaded under ordinary task pressure~\cite{id260930217}, and OOD detection alone
does not close the failure envelope \cite{id221116158}.
\end{description}

\subsection{Three determinism terms: task, model, serving path}
\label{sec:determinism}
\cref{sec:selection} routes tasks on two axes: how open the answer space is, and what an error
costs. That routing requires one distinction we fix here. A \emph{task} has a well-defined or an
open-ended answer space; a \emph{model} either samples stochastically or is effectively
deterministic at inference; and a \emph{serving path} either reproduces the same arithmetic for
the same request or does not, because the inference engine, the numeric precision and the
co-resident traffic on the device all enter the computed function.
The first two are frequently conflated in vendor material (``deterministic mode'' versus
``deterministic task''), and we keep them strictly separated. The third is the one a deployment
team discovers by accident: no checkpoint records it, so a threshold fitted on one serving path
does not carry to another (\cref{sec:l5-bridge}). The planner on a well-posed route
query and the generative assistant on an ambiguous brief face different reliability demands
even under identical decoding parameters.

\subsection{Vocabulary used in this paper}
\begin{description}[leftmargin=1.5em, itemsep=1pt]
\item[Layer.] One of the reliability-bearing decision problems the stack arranges, from task and
risk specification to assurance and control; a decision problem, not a model class.
\item[Stack instance.] A particular filling of the layers by one assembled system; the unit is
not a model considered on its own.
\item[Inherited failure envelope.] The failure modes a layer's choice makes unrecoverable
downstream: the choice cannot remove them, only relocate them.
\item[Reliability debt.] What a fix incurs when it moves a failure to another layer rather than
removing it; see also \emph{inherited failure envelope}.
\item[Well-posedness.] A task is well-posed when its answer space and input distribution can both
be stated; reliability is claimed only relative to that pair.
\item[Selective prediction.] Turning measured confidence into a decision by fixing a cutoff with a
stated error rate.
\item[Validity versus truth.] Schema or syntactic correctness of an output, which a contract layer
can enforce, versus semantic correctness of its content, which no contract can.
\item[Residual.] The error left after the safeguards act, which the operator chooses to own and
re-evaluate rather than assume absent.
\end{description}
